\scanpage{4}
Si $E$ est muni d'une base
\[
 (e_1,e_2,e_3)\ \unsure{\in|e}\qquad (e_i\in\Gamma(S,E)) \tag*{(8)}
\]
définissant la base duale
\[
 (e'_1,e'_2,e'_3)\qquad(e'_i\in\Gamma(S,E')) \tag*{(8$'$)}
\]
par
\[
 \langle e_i,e'_j\rangle=\delta_{ij}\qquad\text{(symb. de Kronecker)}
\]
on a des bases de $E,E'$ par
\[
\left\{\begin{aligned}
 \bar e_1=e_2\wedge e_3,\ \bar e_2=e_3\wedge e_1,\ \bar e_3=e_1\wedge e_2
 &\qquad(\text{base de }\bigwedge^2 E)\\
 \bar e'_1=e'_2\wedge e'_3,\ \bar e'_2=e'_3\wedge e'_1,\ \bar e'_3=e'_1\wedge e'_2
 &\qquad(\text{--- }\bigwedge^2 E')
\end{aligned}\right. \tag*{(\unsure{9|8})}
\]
et on a
\[
 \left\{\begin{aligned}
 \alpha(\bar e_i)&=e'_i\\
 \alpha'(\bar e'_i)&=e_i
 \end{aligned}\right. . \tag*{(\unsure{10|9})}
\]
Notons \deleted{\unsure{donc}}

\uline{Observation} Si $E$ est un Module vectoriel ``spécial'' de rang 3, alors on a \deleted{une} bijection\deleted{\unsure{s}} canoniques
\diagram{\begin{tikzcd}[ampersand replacement=\&,column sep=small]
{\substack{\text{Formes bilinéaires}\\ E\times E\to\underline{\mathrm O}_S}}
 \arrow[rr,"\sim"] \arrow[dr,"\sim"] \& \&
 {\substack{\text{structures d'Algèbre}\\ \text{alternée sur } E'}}\\
 \unsure{\text{---}} \& \operatorname{Hom}_{\underline{\mathrm O}_S}(E,E')\arrow[ur,"\sim"] \&
\end{tikzcd}}
où $\overset{\sim}{\searrow}$ \uline{associe} à $\underline{\varphi(x,y)}$ \uline{l'application $f$}$:E\to E'$ définie par
\[
 f(x)(y)\bigl(=\langle y,f(x)\rangle\bigr)=\varphi(x,y) \tag*{(11)}
\]
et où $\overset{\sim}{\nearrow}$ est l'application qui à $\deleted{\unsure{\varphi}}f$ associe la structure d'algèbre alternée
\[
 (x',y')\longmapsto [x',y']_\varphi=\deleted{\unsure{\varphi}}f\alpha'(x'\wedge y')\ . \tag*{(12)}
\]

% archive sheet number: 3
